% TODO make hw stylesheet

\documentclass[10pt]{article}
\usepackage[letterpaper,top=1in,left=1in,right=1in]{geometry}
\usepackage[dvipsnames,svgnames]{xcolor}
\usepackage[shortlabels]{enumitem}
\usepackage[framemethod=TikZ]{mdframed}
\usepackage{amsmath,amssymb,amsthm}
\usepackage[colorlinks]{hyperref}
\usepackage{multicol}
\usepackage{tabularx}
\usepackage{bbm}
\usepackage{tikz-cd}

\hypersetup{
    colorlinks=true,
    linkcolor=blue,
    filecolor=magenta,
    urlcolor=magenta,
    pdftitle={MAT 535 HW}
}

\usepackage{mathtools}
\usepackage{thmtools}
\usepackage[textsize=footnotesize]{todonotes}
\usepackage{titling}
\usepackage{cancel}

\reversemarginpar
\mdfdefinestyle{mdbluebox}{roundcorner=10pt,innerbottommargin=9pt,
linecolor=blue,backgroundcolor=TealBlue!5,}
\declaretheoremstyle[headfont=\sffamily\bfseries\color{MidnightBlue},
mdframed={style=mdbluebox},]{thmbluebox}

\mdfdefinestyle{mdbloobox}{frametitlefont=\bfseries,innerbottommargin=8pt,
nobreak=true,backgroundcolor=TealBlue!5,linecolor=blue,}
\declaretheoremstyle[headfont=\bfseries\color{MidnightBlue},
mdframed={style=mdbloobox},headpunct={\\[3pt]},postheadspace=0pt,]{thmbloobox}

\mdfdefinestyle{mdredbox}{frametitlefont=\bfseries,innerbottommargin=8pt,
nobreak=true,backgroundcolor=Salmon!5,linecolor=RawSienna,}
\declaretheoremstyle[headfont=\bfseries\color{RawSienna},
mdframed={style=mdredbox},headpunct={\\[3pt]},postheadspace=0pt,]{thmredbox}

\mdfdefinestyle{mdgreenbox}{linecolor=ForestGreen,backgroundcolor=ForestGreen!5,
linewidth=2pt,rightline=false,leftline=true,topline=false,bottomline=false,}
\declaretheoremstyle[headfont=\bfseries\sffamily\color{ForestGreen!70!black},
mdframed={style=mdgreenbox},headpunct={ --- },]{thmgreenbox}

\mdfdefinestyle{mdorangebox}{linecolor=RedOrange,backgroundcolor=RedOrange!5,
linewidth=2pt,rightline=false,leftline=true,topline=false,bottomline=false,}
\declaretheoremstyle[headfont=\bfseries\sffamily\color{RedOrange!70!black},
mdframed={style=mdorangebox},headpunct={ --- },]{thmorangebox}

\mdfdefinestyle{mdblackbox}{linecolor=black,backgroundcolor=RedViolet!5!gray!5,
linewidth=3pt,rightline=false,leftline=true,topline=false,bottomline=false,}
\declaretheoremstyle[mdframed={style=mdblackbox}]{thmblackbox}

\declaretheoremstyle[spaceabove=6pt,spacebelow=6pt]{boldhead}
\declaretheoremstyle[spaceabove=6pt,spacebelow=6pt,headfont=\normalfont\itshape]{italichead}

\declaretheorem[style=boldhead,name=Problem]{problem}
\declaretheorem[style=boldhead,name=Problem,numbered=no]{problem*}
\declaretheorem[style=boldhead,name=Setup]{setup} % for config geo
\declaretheorem[style=boldhead,name=Example,sibling=problem]{example}
\declaretheorem[style=boldhead,name=$\bigstar$ Problem,sibling=problem]{niceprob}
\declaretheorem[style=boldhead,name=Theorem,sibling=problem]{theorem}
\declaretheorem[style=boldhead,name=Corollary,sibling=theorem]{corollary}
\declaretheorem[style=boldhead,name=Proposition,sibling=theorem]{prop}
\declaretheorem[style=boldhead,name=Lemma,numberwithin=problem]{lemma}
\declaretheorem[style=boldhead,name=Theorem,numbered=no]{theorem*}
\declaretheorem[style=boldhead,name=Lemma,numbered=no]{lemma*}
\declaretheorem[style=boldhead,name=Claim,numberwithin=problem]{claim}
\declaretheorem[style=boldhead,name=Claim,numbered=no]{claim*}
\declaretheorem[style=boldhead,name=Remark,numberwithin=problem]{remark}
\declaretheorem[style=boldhead,name=Remark,numbered=no]{remark*}
\declaretheorem[style=boldhead,name=Definition,sibling=theorem]{definition}
\declaretheorem[style=boldhead,name=Definition,numbered=no]{definition*}

\providecommand{\ol}{\overline}
\providecommand{\tensor}{\otimes}
\renewcommand{\hom}{\operatorname{Hom}}
\providecommand{\tor}{\operatorname{Tor}}
\providecommand{\Gal}{\operatorname{Gal}}
\providecommand{\ceil}[1]{\left\lceil #1 \right\rceil}
\providecommand{\floor}[1]{\left\lfloor #1 \right\rfloor}
\newcommand{\dd}{\,\mathrm{d}}
\providecommand{\ul}{\underline}
\providecommand{\eps}{\varepsilon}
\providecommand{\half}{\frac{1}{2}}
\providecommand{\CC}{\mathbb C}
\providecommand{\FF}{\mathbb F}
\providecommand{\II}{\mathbb I}
\providecommand{\NN}{\mathbb N}
\providecommand{\QQ}{\mathbb Q}
\providecommand{\RR}{\mathbb R}
\providecommand{\ZZ}{\mathbb Z}
\providecommand{\ind}{\mathbbm 1}
\providecommand{\dg}{^\circ}
\providecommand{\ii}{\item}
\providecommand{\setdiff}{\smallsetminus}
\providecommand{\alert}[1]{{\sffamily\textbf{\textcolor{blue}{#1}}}}
\providecommand{\inv}{^{-1}}
\providecommand{\mb}[1]{\mathbf{#1}}
\newcommand{\what}{\widehat}

\newenvironment{soln}{\begin{proof}[Solution]}{\end{proof}}

\providecommand{\pow}{\operatorname{Pow}}
\providecommand{\vocab}[1]{{\textbf{\textcolor{ForestGreen}{#1}}}}
\providecommand{\lcm}{\operatorname{lcm}}
\providecommand{\abs}[1]{\left\lvert #1\right\rvert}
\providecommand{\smabs}[1]{\lvert #1\rvert}
\providecommand{\innerprod}[1]{\langle\!\langle #1\rangle\!\rangle}
\providecommand{\norm}[1]{\left\| #1\right\|}
\providecommand{\smnorm}[1]{\| #1\|}
\providecommand{\gr}{\operatorname{gr}}

\usepackage{mathrsfs}

% place title at top right
\setlength{\droptitle}{-8ex}
\pretitle{\begin{flushright}\Large\bfseries}
\posttitle{\par\end{flushright}}
\preauthor{\begin{flushright}}
\postauthor{\end{flushright}}
\predate{\begin{flushright}}
\postdate{\end{flushright}}

\title{MAT 535 HW06}
\author{\large Aditya Pahuja}
\date{\large Due 03/26}
\begin{document}
\maketitle

\begin{problem}
    [\S14.2\#5]
    Prove that the Galois group of $x^p - 2$ is isomorphic
    to the group of matrices
    \begin{equation*}
	G = \left\{\begin{pmatrix}
	    a&b\\0&1
	\end{pmatrix} : a,b\in\FF_p,\; a\ne 0\right\}.
    \end{equation*}
\end{problem}

\begin{soln}
    Observe that the splitting field of $x^p - 2$
    must contain $\sqrt[p]2$,
    whose minimal polynomial has degree $p$, and
    a primitive $p$th root of unity $\zeta$,
    whose minimal polynomial has degree $p-1$.
    Since $x^p - 2$ splits in $\QQ(\sqrt[p]2,\zeta)$,
    this extension must be a splitting field for $x^p-2$,
    and $[\QQ(\sqrt[p]2,\zeta) : \QQ] = p(p-1)$ because
    the degree is bounded above by $p(p-1)$
    and divisible by $p$ and $p-1$.

    The Galois group of this splitting extension is generated
    by the $\QQ$-automorphisms $\rho$ and $\sigma$ defined by
    \begin{gather*}
	\rho(\sqrt[p]2) = \zeta\sqrt[p]2,\quad \rho(\zeta) = \zeta,\\
	\sigma(\sqrt[p]2) = \sqrt[p]2,\quad \sigma(\zeta) = \zeta^g
    \end{gather*}
    where $g$ is chosen so that $\zeta\mapsto\zeta^g$
    generates the cyclic group $\Gal(\QQ(\zeta)/\QQ)\cong(\ZZ/p\ZZ)^\times$,
    or equivalently $g$ generates $(\ZZ/p\ZZ)^\times$.
    The defining relations between $\sigma$ and $\rho$ are
    \begin{equation*}
	\rho^p = \sigma^{p-1} = 1,\qquad \sigma\rho\sigma\inv = \rho^g
    \end{equation*}
    (these determine the Galois group uniquely
    because the Sylow theorems imply that the subgroup of order $p$ is normal,
    so that the Galois group is a semidirect product
    $\ZZ/p\ZZ \rtimes \ZZ/(p-1)\ZZ$).

    Consider the map $\Phi\colon\Gal(\QQ(\sqrt[p]2,\zeta)/\QQ)\to G$ sending
    \begin{equation*}
        \Phi(\rho) = \begin{pmatrix}
	    1&1\\0&1
        \end{pmatrix},\qquad
	\Phi(\sigma) = \begin{pmatrix}
	    g&0\\0&1
	\end{pmatrix}.
    \end{equation*}
    To check that this is an isomorphism
    we just need to check that $\Phi$ commutes with
    the defining relations, since this implies that
    $G$ has the same presentation as the Galois group.
    \begin{equation*}
        \Phi(\rho^k) = \begin{pmatrix}
	    1&k\\0&1
        \end{pmatrix}
    \end{equation*}
    for any integer $k$ implies that $\Phi(\rho)$ has order $p$;
    \begin{equation*}
        \Phi(\sigma^k) = \begin{pmatrix}
	    g^k & 0\\0 & 1
        \end{pmatrix}
    \end{equation*}
    implies that $\Phi(\sigma)$ has order $p-1$; and
    \begin{equation*}
        \Phi(\sigma\rho\sigma\inv)
	= \begin{pmatrix}
	    g&0\\0&1
	\end{pmatrix}
	\begin{pmatrix}
	    1&1\\0&1
	\end{pmatrix}
	\begin{pmatrix}
	    g\inv&0\\0&1
	\end{pmatrix}
	= \begin{pmatrix}
	    1&g\\0&1
	\end{pmatrix} = \Phi(\rho^g).\qedhere
    \end{equation*}
\end{soln}

\pagebreak
\begin{problem}
    [\S14.2\#10]
    Determine the Galois group over $\QQ$ of $x^8 - 3$.
\end{problem}

\begin{soln}
    The splitting field of $x^8 - 3$ is the degree $32$ extension
    $E = \QQ(\sqrt[8]3, \zeta)$ where
    $\zeta = \frac{\sqrt 2}2 + i\frac{\sqrt 2}2$
    is a primitive $8$th root of unity; the extension has degree $32$
    because $[E : \QQ(\sqrt[8]3)] = 4$,
    since the smaller field only contains real numbers.

    Since $E/\QQ(\sqrt[8]3)$ is a biquadratic extension
    generated by $i$ and $\sqrt 2$,
    \[
	\Gal(E/\QQ(\sqrt[8]3))\cong V_4 = \ZZ/2\ZZ\times\ZZ/2\ZZ,
    \]
    consisting of the automorphisms
    $\sigma$ given by complex conjugation
    (hence swapping $\zeta\leftrightarrow\zeta^7$
    and $\zeta^3\leftrightarrow \zeta^5$);
    $\tau$ given by negating $\zeta$
    (hence swapping $\zeta\leftrightarrow\zeta^5$
    and $\zeta^3\leftrightarrow\zeta^7$);
    and $\sigma\tau$ (which swaps $\zeta\leftrightarrow\zeta^3$
    and $\zeta^5\leftrightarrow\zeta^7$).

    We also see that $\Gal(E/\QQ(\zeta))$ contains the automorphism
    $\varphi$ sending $\sqrt[8]3\mapsto\zeta\cdot \sqrt[8]3$.
    The powers of $\varphi$ send
    $\varphi^k(\sqrt[8]3) = \zeta^k\sqrt[8]3$, so
    $\varphi$ has order $8$, meaning it generates
    the order $8$ group $\Gal(E/\QQ(\zeta))\cong\ZZ/8\ZZ$.
    Additionally, $\QQ(\zeta)/\QQ$ is a Galois extension
    (since $\QQ(\zeta)$ is the splitting field of $x^4 + 1$),
    so $\Gal(E/\QQ(\zeta))$ is normal in $\Gal(E/\QQ)$.

    Finally, observe that $\Gal(E/\QQ(\zeta))\cap\Gal(E/\QQ(\sqrt[8]3))$
    fixes the compositum $\QQ(\zeta)\QQ(\sqrt[8]3) = E$,
    so the intersection of these subgroups is the trivial subgroup,
    while $\Gal(E/\QQ(\zeta))\cdot\Gal(E/\QQ(\sqrt[8]3))$ fixes
    $\QQ(\zeta)\cap\QQ(\sqrt[8]3) = \QQ$, hence equals $\Gal(E/\QQ)$.
    This means $\Gal(E/\QQ)\cong \ZZ/8\ZZ\rtimes_\theta V_4$,
    with the action $\theta$ of $V_4$ on $\ZZ/8\ZZ$ given by
    \begin{align*}
        \theta_\sigma(\varphi) = \sigma\varphi\sigma
	&=
	\begin{cases}
	    \zeta&\mapsto\zeta\\
	    \sqrt[8]3&\mapsto\zeta^7\cdot\sqrt[8]3
	\end{cases}\\
	\theta_\tau(\varphi) = \tau\varphi\tau
	&=
	\begin{cases}
	    \zeta&\mapsto\zeta\\
	    \sqrt[8]3&\mapsto\zeta^5\cdot\sqrt[8]3
	\end{cases}\\
	\theta_{\sigma\tau}(\varphi) = (\sigma\tau)\varphi(\sigma\tau)
	&=
	\begin{cases}
	    \zeta&\mapsto\zeta\\
	    \sqrt[8]3&\mapsto\zeta^3\cdot\sqrt[8]3,
	\end{cases}
    \end{align*}
    or equivalently,
    \begin{equation*}
        \Gal(E/Q) = \langle\sigma,\tau,\varphi\mid
	\varphi^8 = \sigma^2 = \tau^2 = 1,\;
	\sigma\varphi\sigma = \varphi^7,\;
	\tau\varphi\tau = \varphi^5\rangle
    \end{equation*}
    (noting that this implies
    $(\sigma\tau)\varphi(\sigma\tau)
    = \tau(\sigma\varphi\sigma)\tau
    = \tau\varphi^7\tau
    = \varphi^{35} = \varphi^3$).
\end{soln}

\pagebreak
\begin{problem}
    [\S14.2\#16]
    \begin{enumerate}[(a)]
        \ii Prove that $x^4 - 2x^2 - 2$ is irreducible over $\QQ$.
	\ii Show that the roots of this quartic are
	\begin{equation*}
	    \alpha_1 = \sqrt{1 + \sqrt 3},\quad
	    \alpha_2 = \sqrt{1 - \sqrt 3},\quad
	    \alpha_3 = -\sqrt{1 + \sqrt 3},\quad
	    \alpha_4 = -\sqrt{1 - \sqrt 3}.
	\end{equation*}
	\ii Let $K_1 = \QQ(\alpha_1)$ and $K_2 = \QQ(\alpha_2)$.
	Show that $K_1\ne K_2$ and $K_1\cap K_2 = \QQ(\sqrt 3) = F$.
	\ii Show that $K_1$, $K_2$, and $K_1K_2$ are Galois over $F$
	with $\Gal(K_1K_2/F)$ the Klein $4$-group.
	Write out the elements of $\Gal(K_1K_2/F)$ explicitly.
	Determine all the subgroups of the Galois group
	and give their corresponding fixed subfields of
	$K_1K_2$ containing $F$.
	\ii Prove that the splitting field of $x^4 - 2x^2 - 2$ over $\QQ$
	is of degree $8$ with dihedral Galois group.
    \end{enumerate}
\end{problem}

\begin{soln}
    (a) Follows from Eisenstein's criterion with the prime $2$.
    \\

    (b) Follows from writing
    $x^4 - 2x^2 - 2 = (x^2 - 1)^2 - 3$,
    since setting this equal to zero implies
    $x^2 - 1 = \pm 3$.
    \\

    (c) On one hand, $K_1\cap K_2$ contains
    $\alpha_1^2 - 1 = 1 - \alpha_2^2 = \sqrt 3$,
    so it contains $\QQ(\sqrt 3)$.
    On the other hand, $K_1\cap K_2$ is a proper subextension of $K_1/\QQ$,
    so its degree over $\QQ$ is a proper divisor of $[K_1 : \QQ] = 4$,
    meaning the degree is at most $2$. This shows that
    $K_1\cap K_2$ cannot be larger than $\QQ(\sqrt 3)$.
    \\

    (d) We can see that $K_1K_2 = \QQ(\alpha_1,\alpha_2)$
    is a splitting field of $x^4 - 2x^2 - 2$, so certainly
    $K_1K_2/F$ is a Galois extension. Since
    \[
	2 = [K_1 : F] < [K_1K_2 : F] \le [K_1 : F][K_2 : F] = 4
    \]
    and $[K_1K_2 : F]$ is divisble by $[K_1 : F] = 2$,
    the degree of $K_1K_2/F$ is $4$.
    Furthermore $\Gal(K_1K_2/F)$ contains the two order-$2$ automorphisms
    $\sigma$ and $\tau$ defined by
    \begin{gather*}
        \sigma(\alpha_1) = \alpha_3,\quad
	\sigma(\alpha_2) = \alpha_2,\quad
	\sigma(\alpha_3) = \alpha_1,\quad
	\sigma(\alpha_4) = \alpha_4,\\
        \tau(\alpha_1) = \alpha_1,\quad
	\tau(\alpha_2) = \alpha_4,\quad
	\tau(\alpha_3) = \alpha_3,\quad
	\tau(\alpha_4) = \alpha_2.
    \end{gather*}
    To check that $\sigma$ is an automorphism, we note that
    $\sigma$ fixes $K_2$ and swaps the roots of the polynomial
    $x^2 - 1 - \sqrt 3$, hence factors through the map
    $K_2[x]\to K_2(\alpha_1) = K_1K_2$ sending $x$ to $\alpha_1$,
    so that $\sigma$ is a field homomorphism (hence isomorphism,
    by degree considerations)
    $K_2[x]/(x^2 - 1 - \sqrt3)\cong K_1K_2\to K_1K_2$;
    similar reasoning shows that $\tau$ is a field isomorphism.

    The only groups of order $4$ up to isomorphism
    are $\ZZ/4\ZZ$ and $\ZZ/2\ZZ\times\ZZ/2\ZZ$,
    but the former only has one element of order $2$,
    so $\Gal(K_1K_2/F)\cong\ZZ/2\ZZ\times\ZZ/2\ZZ$
    (and the remaining non-identity automorphism is $\sigma\tau=\tau\sigma$,
    which swaps both pairs $\{\alpha_1,\alpha_3\}$
    and $\{\alpha_2,\alpha_4\}$).

    Since $G = \Gal(K_1K_2/F)$ is abelian,
    its subgroups are all normal subgroups of $G$,
    meaning the subextensions $K_1/F$ and $K_2/F$ of $K_1K_2/F$
    are both Galois extensions, by the fundamental theorem.
    Furthermore, the subgroups of $G$ are
    $\{1\}$, $\{1,\sigma\}$, $\{1,\tau\}$, $\{1,\sigma\tau\}$, and $G$,
    which respectively have fixed fields
    $K_1K_2$, $K_2$, $K_1$, $F(\alpha_1\alpha_2) = F(\sqrt{-2})$, and $F$.
    \\

    (e) The degree of the splitting field $E = K_1K_2$ of $x^4 - 2x^2 - 2$
    over $\QQ$ is $[E : F][F : \QQ] = 4\cdot 2 = 8$.
    Consider $\Gal(E/\QQ)$ as a subgroup of $S_4$,
    with $k=1,2,3,4$ corresponding to $\alpha_k$
    (so $\sigma$ from part (d) is the permutation $(1\;3)$).
    Then the permutations in (d) are
    \begin{equation*}
        1 = (1),\quad
	\sigma = (1\;3),\quad
	\tau = (2\;4),\quad
	\sigma\tau = \tau\sigma = (1\;3)(2\;4).
    \end{equation*}
    The remaining four automorphisms are the ones not fixing $F$;
    one such automorphism is $(1\;2)(3\;4)$.
    This is an automorphism because it fixes the field
    $\QQ(\beta)$ where $\beta^2 = 2 + 2\sqrt{-2}$
    and swaps the roots of $x^2 - \beta x + \sqrt{-2}
    = (x-\alpha_1)(x-\alpha_2)$ (similar to the logic for $\sigma$
    being an automorphism in (d)). The coset of $\Gal(E/F)$
    by (the automorphism corresponding to) $(1\; 2)(3\; 4)$
    is $\{(1\;2)(3\;4),(1\;2\;3\;4),(1\;4\;3\;2),(1\;4)(2\;3)\}$.
    Thus, $\Gal(E/\QQ)$ is generated as a subgroup of $S_4$ by
    $\varphi = (1\;2\;3\;4)$ and $\sigma$, which satisfy
    $\varphi^4 = \sigma^2 = (\sigma\varphi)^2 = 1$, meaning it is dihedral.
\end{soln}

\begin{problem}
    [14.4\#5]
    Let $p$ be a prime and let $F$ be a field.
    Let $K$ be a Galois extension of $F$ whose Galois group is a $p$-group.
    Such an extension is called a $p$-extension.
    \begin{enumerate}[(a)]
        \ii Let $L$ be a $p$-extension of $K$.
	Prove that the Galois closure of $L$ over $F$
	is a $p$-extension of $F$.
	\ii Give an example to show that (a) need not hold
	if $[K : F]$ is a power of $p$ but $K/F$ is not Galois.
    \end{enumerate}
\end{problem}

\begin{soln}
    (a)
    Let $\ol F$ be the Galois closure of $L$ over $F$
    and define $G = \Gal(\ol F/F)$,
    $N = \Gal(\ol F/K)$, and $Q = \Gal(\ol F/L)$.
    The given setup translates to the following
    group-theoretic information:
    \begin{itemize}
	\ii $N\trianglelefteq G$ and $Q\trianglelefteq N$,
	since $K/F$ and $L/K$ are Galois;
	\ii $(G:N) = [K:F]$ and $(N:Q) = [L:K]$ are both powers of $p$; and
	\ii $Q$ does not contain any nontrivial normal subgroups of $G$,
	since no proper $L$-subextension of $\ol F$ is Galois over $F$.
    \end{itemize}
    Our goal is to show that $G$ (and $Q$ and $N$) are $p$-groups.

    Let $Q = Q_1$, $Q_2$, \dots, $Q_n$ be the distinct
    conjugate subgroups of $Q$ in $G$, so
    \[
	\{Q_1,\dots,Q_n\} = \{gQg\inv : g\in G\}.
    \]
    Since $N$ is normal in $G$ and $Q$ is normal in $N$,
    the conjugates of $Q$ are each normal subgroups of $N$.
    Furthermore, since $Q$ does not contain any nontrivial subgroups of $G$
    and the intersection of the $Q_i$ is by definition
    normal in $G$, we find that $Q_1\cap\cdots\cap Q_n = \{1\}$.
    Then, recalling that in general
    $\abs{H\cdot K} = \frac{\abs H\abs K}{\abs{H\cap K}}$
    for subgroups $H,K$ of a finite group, we can compute
    \begin{equation*}
	(G : Q_1\cap\cdots\cap Q_k)
	= \frac{(G : Q_1\cap\cdots\cap Q_{k-1})(G : Q_k)}
	{(G:(Q_1\cap\cdots\cap Q_{k-1})\cdot Q_k)},
    \end{equation*}
    so in particular $(G:Q_1\cap\cdots\cap Q_k)$
    divides $(G:Q_1\cap\cdots\cap Q_{k-1})(G:Q_k)$.
    Since $(G:Q_i)$ is a power of $p$ for each $i=1,\dots,n$,
    we have shown by induction that $(G:Q_1\cap\cdots\cap Q_k)$
    is a (factor of a) power of $p$ for each $k=1,\dots,n$,
    so $(G:Q_1\cap\cdots\cap Q_n) = (G:1) = \abs G$ is a power of $p$.
    \\

    (b) For odd primes $p$, let $F = \QQ$ and $K = L = \QQ(\sqrt[p]2)$.
    We can see that $[K : F] = p$ and $K/F$ is not a Galois extension
    (since $x^p - 2$ has a root in $K$ but doesn't split),
    and $L/K$ is trivially a $p$-extension.
    Then the Galois closure of $L$ over $\QQ$ is a subextension of
    $E = L(\zeta)$ which strictly contains $L$,
    where $\zeta$ is a root of $1 + x + \cdots + x^{p-1}$,
    since $\QQ(\sqrt[p]2,\zeta)$, being the splitting field of $x^p - 2$,
    is Galois over $\QQ$. Since the degree of $\ol F$ in $L$
    divides $[E : L] = p-1$ and is not equal to $1$,
    $\ol F$ cannot be a $p$-extension.
\end{soln}

\pagebreak

\begin{problem}
    [\S14.4\#8]
    Let $K_1$ and $K_2$ be two algebraic extensions of a field $F$
    contained in the field $L$ of characteristic zero.
    Prove that the $F$-algebra $K_1\tensor_F K_2$
    has no nonzero nilpotent elements.
\end{problem}

\begin{soln}
    Consider an arbitrary element
    $\sum_{i=1}^n \alpha_i\tensor\beta_i\in K_1\tensor_F K_2$
    for $\alpha_i\in K_1$ and $\beta_i\in K_2$.
    Then in fact $c$ lies in the sub-$F$-algebra
    $K_1\tensor_F F(\beta_1,\dots,\beta_n)$.
    Since $L$ has characteristic zero, $F(\beta_1,\dots,\beta_n)$
    is a separable, finite extension, so it is a simple extension
    $F(\theta)$ by the primitive element theorem.

    \begin{lemma}
	[\S14.4\#7]
        Let $F\subseteq K\subseteq L$ and let $\theta\in L$
	with $p(x) = m_{\theta,F}(x)$.
	Prove that $K\tensor_F F(\theta) \cong K[x]/(p(x))$ as $K$-algebras.
    \end{lemma}

    \begin{proof}
        Let $d = \deg p$, so $F(\theta)$ has $F$-basis
	$1$, $\theta$, \dots, $\theta^{d-1}$. Then
	consider the map $\varphi\colon K[x]\to K\tensor_F F(\theta)$
	defined by
	\begin{equation*}
	    \varphi(a_0 + a_1x + \dots + a_mx^m)
	    = a_0\tensor 1 + a_1\tensor\theta
	    + \cdots + a_m\tensor \theta^m.
	\end{equation*}
	To verify that this is a $K$-algebra homomorphism
	we note that $\varphi(1) = 1\tensor 1$ is the identity,
	\begin{equation*}
	    \varphi(ax^p\cdot bx^q) = ab\tensor x^{p+q}
	    = (a\tensor x^p)\cdot (b\tensor x^q)
	\end{equation*}
	so that $\varphi$ commutes with multiplication of monomials,
	and then it commutes with multiplication of arbitrary polynomials
	using $K$-linearity. Additionally, $\varphi$ is clearly surjective.

	The kernel of $\varphi$ contains $p(x)$ because
	\begin{equation*}
	    \varphi(p(x)) = 1\tensor p(\theta) = 0,
	\end{equation*}
	noting that $p(x)\in F[x]$ means that
	we can move the coefficients of $p$ through the $\tensor$.
	On the other hand, since $F(\theta)$ is isomorphic to
	$F^d$ as an $F$-vector space, $K\otimes_F F(\theta)\cong K^d$
	as a $K$-vector space, which means that
	$K[x]/\ker\varphi$ must be a $d$-dimensional $K$-vector space.
	This forces the generator of $\ker\varphi$
	to be a polynomial of degree $d$, so $\ker\varphi = (p(x))$,
	implying $K[x]/(p(x))\cong K\otimes_F F(\theta)$ as desired.
    \end{proof}

    The lemma implies that there is a $K$-algebra isomorphism
    between $K_1\otimes_F F(\theta)$ and $K[x]/(p(x))$.
    The latter has no nonzero nilpotent elements (because it is a field),
    so the former also has no nonzero nilpotent elements;
    thus the arbitrarily chosen element $\sum \alpha_i\otimes\beta_i$
    cannot be nilpotent unless it is zero.
\end{soln}










\end{document}
